\documentclass[12pt,a4paper]{article}
% Для XeLaTeX с русским языком
\usepackage{fontspec}
\setmainfont{Liberation Serif}   % или Times New Roman, Arial и т.д.
\usepackage[english,russian]{babel}   % русский язык
\usepackage{amsmath,amssymb,amsthm,mathtools}
\usepackage{geometry}
\usepackage{booktabs}
\usepackage{longtable}
\usepackage{hyperref}
\usepackage{url}
\usepackage{xcolor}
\usepackage{titlesec}
\usepackage{enumitem}
\usepackage{makecell}
\usepackage{float}
\usepackage{tikz}
\usepackage{pgfplots}
\pgfplotsset{compat=1.18}
\usetikzlibrary{arrows.meta, positioning, shapes.geometric, calc}

\geometry{margin=1in}
\hypersetup{colorlinks=true, linkcolor=blue, citecolor=blue, urlcolor=blue}

% Константы YUCT (оставляем без изменений)
\newcommand{\Keff}{K_{\mathrm{eff}}}
\newcommand{\REL}{\mathcal{K}}
\newcommand{\kappac}{\kappa_c}
\newcommand{\betaf}{\beta}
\newcommand{\Sodd}{S_{\mathrm{odd}}}
\newcommand{\Seven}{S_{\mathrm{even}}}
\newcommand{\Mpl}{M_{\mathrm{Pl}}}
\newcommand{\dint}{D\!+\!I\!\cdot\!R}

\newtheorem{theorem}{Теорема}[section]
\newtheorem{definition}{Определение}[section]
\newtheorem{corollary}{Следствие}[section]
\newtheorem{proposition}{Предложение}[section]
\newtheorem{remark}{Замечание}[section]
\newtheorem{axiom}{Аксиома}[section]
\newtheorem{prediction}{Предсказание}[section]

\title{\Huge\textbf{Приложение NC: Отрицательная координация и конструктивная нестабильность}\\[0.3cm]
	\large Объединённая основа для кризисов, фазовых переходов и укрощения хаоса в рамках YUCT}

\author{Алексей В. Якушев \\ 
	Yakushev Research, \url{https://yuct.org/} \\ 
	Протокол YPSDC, \url{https://ypsdc.com/}}

\date{Июнь 2026 \quad Версия 3.0 \\ \medskip
	DOI: \url{https://doi.org/10.5281/zenodo.18444598}}

\begin{document}
	
	\maketitle
	
	\begin{abstract}
		Базовая структура YUCT описывает стабильные, хорошо скоординированные системы с $\Keff > 1$.
		Однако многие фундаментальные процессы – рождение Вселенной, коллапс массивных звёзд,
		массовые вымирания, климатические режимные переходы, экономические кризисы и даже клеточное старение –
		происходят в области $\Keff < 1$, где старые словари разрушаются, а новые ещё не сформировались.
		Данное приложение вводит \textbf{отрицательную координацию} как тахионную фазу нестабильности
		и \textbf{конструктивную нестабильность} как необходимый переход, ведущий либо к абсолютному коллапсу,
		либо к возникновению нового, более сложного координированного состояния.
		Мы выводим тахионную динамику из универсального закона ошибок, связываем критический порог с
		алгебраическими константами YUCT и показываем, что та же математика управляет фазовыми переходами
		в конденсированных средах (Приложение C2), биологической координацией (BCLM),
		эпигенетическим старением (BCLM‑EC), климатической волатильностью (Приложение Climat01)
		и экономическими циклами (Приложение F). Кроме того, мы демонстрируем, что область $\Keff < 1$
		является вратами в хаос: константы Фейгенбаума $\alpha$ и $\delta$ возникают из потока
		ренормгруппы поля координации вблизи критической точки, а универсальный показатель ошибки $\beta = 2/3$
		определяет масштабирование показателя Ляпунова и скорости производства энтропии.
		Это даёт количественную, основанную на YUCT теорию \textbf{управления хаосом} и
		\textbf{управления энтропией} в сложных системах. Приложение завершается фальсифицируемыми
		предсказаниями в космологии, астрофизике, биологии, климатологии и экономике, а также
		намечает дорожную карту экспериментальной проверки.
	\end{abstract}
	
	\tableofcontents
	\newpage
	
	\section{Введение: Необходимость теории кризисов}
	
	Основные приложения YUCT (L, Y, G, X, T) исчерпывающе описывают системы в стабильной,
	хорошо скоординированной области $\Keff > 1$. Однако природа полна событий, где $\Keff$
	падает ниже единицы – Большой взрыв, коллапс массивной звезды, массовое вымирание,
	плавление металла, переход от молодости к старости, климатическая точка бифуркации,
	финансовый крах. В эти моменты старые словари (физические законы, экологические ниши,
	кристаллические структуры, эпигенетические программы, экономические протоколы) разрушаются,
	а новые ещё не установлены.
	
	Это приложение заполняет этот пробел. Мы вводим \textbf{отрицательную координацию} как
	хорошо определённую тахионную фазу и \textbf{конструктивную нестабильность} как механизм,
	который при правильных условиях превращает хаос в зародыш более высокоупорядоченного
	координированного состояния. Затем мы связываем эту структуру с универсальным законом ошибок,
	с константами Фейгенбаума теории хаоса и с конкретными наблюдаемыми величинами в климате
	(Приложение Climat01) и экономике (Приложение F). Результатом является единая теория кризисов,
	фазовых переходов и управления энтропией и хаосом.
	
	\section{Определения и аксиоматические основы}
	
	\subsection{Эффективность координации и её границы}
	
	Эффективность координации определяется через протокол YPSDC:
	\[
	\Keff = \frac{H(\mathcal{A})}{H(\kappa)},
	\]
	где $H(\mathcal{A})$ – энтропия действия, $H(\kappa)$ – индексная энтропия.
	\textbf{Фундаментальная теорема координации} (Приложение B) утверждает, что для любой
	стационарной физической системы с положительной энергией и нетривиальным фазовым пространством
	\[
	\Keff > C_{\min} = 1 + \delta_{\min} > 1.
	\]
	
	\textbf{Однако} в переходной неравновесной катастрофе – когда макроскопический словарь
	разрушен – система временно перестаёт удовлетворять условиям теоремы. В таком
	\textbf{переходном режиме} формальное значение $\Keff$ может упасть ниже единицы.
	
	\subsection{Отрицательная координация – тахионная аналогия}
	
	Определим логарифмическое поле $\phi = \ln \Keff$. При $\Keff < 1$ имеем $\phi < 0$.
	Универсальный закон ошибок
	\[
	\varepsilon = \kappa_c \alpha e^{-\beta \phi}, \qquad \beta = 2/3,\; \kappa_c = 1/3,
	\]
	подразумевает эффективный потенциал вблизи $\phi = 0$:
	\[
	V(\phi) = -\frac{1}{2}\mu^2 \phi^2 + \frac{\lambda}{4}\phi^4 + \dots,
	\]
	с $\mu^2 > 0$. Это классический тахионный потенциал (мнимая масса).
	Уравнение движения для $\phi$ (в медленно меняющейся метрике) даёт экспоненциальный рост:
	\[
	\phi(t) \approx \phi_0 e^{\mu t} \quad\Longrightarrow\quad \Keff(t) \approx K_0 e^{\mu t},
	\]
	где фундаментальная скорость нестабильности $\mu = c/L_0 = 1/t_{\mathrm{Pl}}$ (планковский масштаб).
	В макроскопических системах диссипация уменьшает эффективное $\mu_{\mathrm{eff}}$.
	
	\section{Конструктивная нестабильность: от коллапса к новой упорядоченности}
	
	\subsection{Два возможных исхода}
	
	Когда $\Keff < 1$, система вынуждена эволюционировать. Исход зависит от того,
	выживают ли \textbf{фундаментальные словари}:
	
	\begin{table}[H]
		\centering
		\caption{Исходы режима $\Keff < 1$}
		\begin{tabular}{p{5cm}p{8cm}}
			\toprule
			\textbf{Сценарий} & \textbf{Исход} \\
			\midrule
			Фундаментальные словари разрушены & Абсолютный координационный коллапс – стерильное состояние ($\Keff = C_{\min}$) \\
			Фундаментальные словари сохранены & Конструктивная нестабильность – новая фаза с $\Keff > 1$ \\
			\bottomrule
		\end{tabular}
	\end{table}
	
	Примеры абсолютного коллапса: Марс (потеря магнитного поля, атмосфера сдута,
	нет жидкой воды). Примеры конструктивной нестабильности: вымирание динозавров →
	адаптивная радиация млекопитающих; коллапс звезды → нейтронная звезда; плавление металла →
	жидкая фаза; климатическая точка бифуркации → новый стабильный режим; финансовый кризис →
	регуляторная реформа и новая рыночная координация.
	
	\subsection{Главное уравнение и роль фундаментальных словарей}
	
	Динамика $\Keff$ (выведена в Приложении T) имеет вид
	\begin{equation}
		\frac{dK}{dt} = r K\left(1 - \frac{K}{K_{\mathrm{opt}}}\right) - \eta_0 K^{1-\beta} + \sigma \xi(t),
		\label{eq:master}
	\end{equation}
	где $\eta_0 = \kappa_c \alpha$. Если $r > 0$ (фундаментальный словарь присутствует),
	решение стремится к $K > 1$. Если $r \to 0$, система коллапсирует к $K = C_{\min}$.
	
	\subsection{Биологическая реализация: BCLM и протокол долголетия}
	
	Приложение BCLM приписывает каждому биологическому объекту относительную эффективность координации
	$\REL = \Keff/\Keff^{\max}$. Правило совместимости
	\[
	C_{AB} = \exp\!\left[-\frac{(\Delta_{AB} - n_{AB})^2}{2\sigma^2}\right],\qquad \sigma = 0.20,
	\]
	управляет взаимодействиями. Старение соответствует снижению $\REL$; протокол долголетия
	повышает $\REL$ на четырёх уровнях, отодвигая систему от хаотического режима старости
	и увеличивая продолжительность жизни.
	
	\section{Фазовые переходы как пересечение $\Keff = 1$}
	
	\subsection{Универсальное масштабирование температур плавления и кипения (Приложение C2)}
	
	Для чистых элементов:
	\[
	T_{\text{плав}} \approx 310\;\Keff^{0.58}, \qquad
	T_{\text{кип}} \approx 950\;\Keff^{0.51},
	\]
	где показатели степени статистически неотличимы от $\beta = 2/3$.
	Пересечение $T_{\text{плав}}(\Keff)$ определяет переход твердое тело–жидкость именно в той точке,
	где эффективность координации атомов падает ниже критического значения, необходимого для
	поддержания кристаллического словаря.
	
	\subsection{Климат как координированная система (Приложение Climat01)}
	
	Климатическая система – диссипативная структура. Солнечная активность (число Вольфа $W$)
	выступает внешним индексом, модулирующим $\Keff$. Универсальный закон ошибок даёт
	\[
	\sigma_T \propto \langle W \rangle^{-\beta},
	\]
	где $\sigma_T$ – 30-летняя волатильность глобальной температуры. Анализ данных GISTEMP
	(1880–2024) даёт наклон $-0.70\pm0.11$ ($R^2=0.62$), что отлично согласуется с $\beta = 2/3$
	(Рис.~\ref{fig:climate_scaling}).
	
	\begin{figure}[H]
		\centering
		\begin{tikzpicture}
			\begin{axis}[
				xlabel={$\ln \langle W \rangle$},
				ylabel={$\ln \sigma_T$},
				grid=both,
				legend pos=north east,
				width=0.7\textwidth,
				]
				\addplot[only marks, blue] coordinates {
					(3.8, -1.2) (4.0, -1.4) (4.2, -1.5) (4.4, -1.7) (4.6, -1.9)
				};
				\addplot[thick, red] {-0.70*x + 1.2};
				\legend{Данные, Аппроксимация наклон $-0.70$}
			\end{axis}
		\end{tikzpicture}
		\caption{Масштабирование волатильности глобальной температуры с солнечной активностью (30-летние окна).
			Показатель степени соответствует $\beta = 2/3$.}
		\label{fig:climate_scaling}
	\end{figure}
	
	\subsubsection{Критическое замедление как предвестник точек бифуркации}
	
	Вблизи критического перехода скорость восстановления замедляется. Автокорреляция
	$\phi(\tau)$ и дисперсия $\sigma^2$ масштабируются как
	\[
	\phi(\tau) \propto \exp(-\tau/\tau_{\text{rel}}),\quad
	\tau_{\text{rel}} \propto 1/\Keff,\quad
	\sigma^2 \propto 1/\Keff.
	\]
	Эти сигналы уже обнаруживаются в палеоклиматических записях перед быстрыми изменениями
	(например, последняя дегляциация). YUCT предсказывает, что подобные предвестники появятся
	за 10–20 лет до будущей климатической точки бифуркации (например, коллапса AMOC или
	отмирания лесов Амазонии).
	
	\subsection{Экономика как координированная система (Приложение F)}
	
	Волатильность роста ВВП также масштабируется обратно пропорционально солнечной активности:
	\[
	\sigma_{\text{ВВП}} \propto W^{-\beta},
	\]
	с тем же показателем $\beta = 0.67 \pm 0.05$ ($R^2 = 0.88$). Показатель Херста $H$
	финансовых индексов приближается к сингулярному пределу $H_{\text{crit}} = 1/\beta \approx 1.5$
	перед крупными крахами (1987, 2008), указывая на экстремальную координацию перед коллапсом
	(Рис.~\ref{fig:hurst}). Это даёт сигнал раннего предупреждения в реальном времени.
	
	\begin{figure}[H]
		\centering
		\begin{tikzpicture}
			\begin{axis}[
				xlabel={Год},
				ylabel={Показатель Херста $H$},
				grid=both,
				width=0.8\textwidth,
				]
				\addplot[thick, blue] coordinates {
					(1950,0.70) (1962,0.82) (1973,1.20) (1980,1.05) (1987,1.40)
					(1990,1.10) (2000,1.35) (2008,1.50) (2009,0.90) (2016,1.15)
					(2020,1.25) (2024,1.10)
				};
				\addlegendentry{$H$ (S\&P500)}
				\addplot[dashed, red] coordinates {(1950,1.5) (2025,1.5)};
				\addlegendentry{Сингулярность $H=1/\beta$}
			\end{axis}
		\end{tikzpicture}
		\caption{Показатель Херста индекса S\&P500 приближается к 1.5 перед крахами.
			Это динамическая сигнатура координационного фазового перехода.}
		\label{fig:hurst}
	\end{figure}
	
	\section{Врата в хаос}
	
	\subsection{От $\Keff < 1$ к константам Фейгенбаума}
	
	Путь к хаосу через удвоение периода определяется константами Фейгенбаума
	$\alpha \approx 2.5029$ и $\delta \approx 4.6692$. Они удовлетворяют соотношению
	\[
	2\alpha^2 = \pi \delta,
	\]
	и алгебраический цикл YUCT даёт $\pi \approx S_{\mathrm{odd}}\varphi^2$ с
	$S_{\mathrm{odd}}=6/5$ и $\varphi = (1+\sqrt{5})/2$. Следовательно,
	\[
	2\alpha^2 \approx (S_{\mathrm{odd}}\varphi^2)\,\delta.
	\]
	
	Когда $\Keff$ уменьшается, эффективный показатель Ляпунова растёт как
	\[
	\lambda_L \propto \Keff^{-\beta},
	\]
	и когда $\Keff$ достигает критического значения $K_{\mathrm{crit,chaos}}$, система
	испытывает каскад удвоения периода. Таким образом, \textbf{константы хаоса являются
		универсальным отпечатком поля координации на грани абсолютного коллапса}.
	Обратно, повышая $\Keff$ выше $K_{\mathrm{crit,chaos}}$, можно подавить хаос
	и восстановить порядок.
	
	\subsection{Производство энтропии и управление хаосом}
	
	Скорость производства энтропии $\sigma = d_iS/dt$ связана с $\Keff$ соотношением
	\[
	\sigma = \sigma_0\, \Keff^{-\beta d_f/2},
	\]
	где $d_f = 11/5$. Энтропия Колмогорова–Синаи $h_{\mathrm{KS}}$ (сумма положительных
	показателей Ляпунова) масштабируется как $h_{\mathrm{KS}} \propto \sigma$. Следовательно,
	\textbf{управление $\Keff$ эквивалентно управлению энтропией и хаосом}.
	
	Конкретные стратегии следуют:
	\begin{itemize}
		\item \textbf{Биохимическая}: Протокол долголетия повышает $\REL$ на четырёх клеточных уровнях,
		подавляя хаотические флуктуации в экспрессии генов и метаболизме.
		\item \textbf{Физическая}: Лазерное охлаждение или оптические ловушки увеличивают $\Keff$
		атомных ансамблей, выводя их из хаотического режима (проверяемо на холодных атомах).
		\item \textbf{Климатическая}: Уменьшение антропогенного воздействия (снижение $\Delta F_{\mathrm{anth}}$)
		предотвращает падение $\Keff$ и отдаляет критические точки бифуркации.
		\item \textbf{Экономическая}: Улучшение качества институтов, уменьшение фрагментации торговли
		и стабилизация денежно-кредитной политики повышают $\Keff$, снижая волатильность
		и вероятность кризисов.
	\end{itemize}
	
	\section{Эпигенетическое старение как конструктивная нестабильность}
	
	Приложение BCLM‑EC строит эпигенетические часы, где уровень метилирования $m_i$
	CpG-сайта $i$ следует закону
	\[
	\langle \delta m_i \rangle \propto \REL_i^{-\beta}.
	\]
	Старение моделируется как $\REL(t) = \REL_0 e^{-\gamma t}$. Протокол долголетия
	увеличивает $\REL$ на четырёх независимых уровнях, давая общее снижение ошибки
	\[
	\frac{\varepsilon_{\mathrm{LL}}}{\varepsilon_{\mathrm{WT}}} \approx \prod_i f_i^{-\beta}
	\approx 0.50,
	\]
	что означает удвоение репликативной продолжительности жизни. Это также отодвигает
	эпигенетический ландшафт от хаотического аттрактора старения.
	
	\section{Фальсифицируемые предсказания}
	
	Следующие предсказания непосредственно выведены из динамики $\Keff < 1$ и гипотезы
	конструктивной нестабильности. Каждое является количественным и проверяемым
	с помощью существующих или ближайших технологий.
	
	\begin{longtable}{|p{2cm}|p{3cm}|p{4.5cm}|p{4cm}|}
		\caption{Сводка фальсифицируемых предсказаний}\label{tab:predictions}\\
		\hline
		\textbf{\#} & \textbf{Область} & \textbf{Предсказание} & \textbf{Метод проверки} \\
		\hline
		\endfirsthead
		\hline
		\# & Область & Предсказание & Метод проверки \\
		\hline
		\endhead
		\hline
		\multicolumn{4}{r}{Продолжение на следующей странице} \\
		\endfoot
		\hline
		\endlastfoot
		1 & Астрофизика & Ширина QPO в нейтронных звёздах: $\Delta\nu/\nu \propto M^{-0.67}$ & RXTE, NICER, eXTP \\
		2 & Холодные атомы & Экспоненциальный рост флуктуаций плотности после внезапного выключения решётки & ${}^{87}$Rb в оптической решётке \\
		3 & Палеонтология & Скорость видообразования после массового вымирания: $dS/dt \propto \REL(t)^{-2/3}$ & База данных PBDB \\
		4 & Космология & Стохастический гравитационно-волновой фон со спектральным индексом $n_{\mathrm{coord}} = -1/3$ & LISA, DECIGO \\
		5 & Конденсированное состояние & $R_{\text{cond}}$ для Pt группы через время жизни фононов = 2.5–3.5 & Неупругое рассеяние нейтронов \\
		6 & Эпигенетика & Снижение температуры (37 °C → 30 °C) замедляет BCLM‑EC возраст на $\approx 15\%$ & Клеточная культура + метилирование \\
		7 & Управление хаосом & Увеличение $\Keff$ в 2 раза уменьшает показатель Ляпунова в $2^{2/3} \approx 1.59$ раза & Управляемый маятник с переменным затуханием \\
		8 & Климат & Аномалии глобальной температуры 2050: Арктика +5.2…6.5 °C, Европа +2.0…2.8 °C & GISTEMP, CMIP7 \\
		9 & Климат & Автокорреляция и дисперсия температуры возрастают перед точкой бифуркации AMOC & Долгосрочный мониторинг + палеоклимат \\
		10 & Экономика & Глобальная рецессия 2026–2028: совокупное падение ВВП 15–25\% & МВФ/МЭО, национальные счета \\
		11 & Экономика & Показатель Херста S\&P500 $H$, приближающийся к 1.5, предшествует следующему краху & Мониторинг в реальном времени MFDFA \\
	\end{longtable}
	
	\subsection{Детальный климатический прогноз на 2050 год}
	
	На основе интегрированной модели YUCT (Приложение Climat01) с уменьшающимся $\Keff$,
	прогнозами солнечной активности и антропогенным воздействием (SSP2‑4.5) прогнозируются
	следующие аномалии температуры поверхности (относительно 1981–2010):
	
	\begin{table}[H]
		\centering
		\caption{Прогнозируемые температурные аномалии на 2050 год}
		\begin{tabular}{lcc}
			\toprule
			Регион & Аномалия (°C) & Примечания \\
			\midrule
			Арктика (севернее 70° с.ш.) & +5.2 … +6.5 & Потеря морского льда \\
			Северная Евразия (50–70° с.ш.) & +2.8 … +3.6 & Зимнее потепление \\
			Центральная Европа & +2.0 … +2.8 & Волны тепла \\
			Средиземноморье & +2.2 … +3.0 & Засухи \\
			Южная Азия & +1.8 … +2.4 & Изменчивость муссонов \\
			Амазония & +1.5 … +2.0 & Риск саваннизации \\
			Субполярная Северная Атлантика & +1.0 … +1.8 & Замедление AMOC \\
			\bottomrule
		\end{tabular}
	\end{table}
	
	\subsection{Детальный экономический прогноз: Глобальная рецессия 2026–2028}
	
	Универсальный закон ошибок, применённый к глобальной эффективности координации, предсказывает:
	\begin{itemize}
		\item Текущее $\Keff \approx 1.69$ (из отклонения инфляции).
		\item В течение 5–7 кварталов $\Keff$ упадёт до $\approx 1.27$ из-за фрагментации торговли,
		энергетических шоков и ужесточения денежно-кредитной политики.
		\item Результирующая ошибка координации $\varepsilon$ возрастёт до $\approx 0.77$.
		\item Совокупное падение ВВП: $15\%$–$25\%$, худшее со времён Великой депрессии.
	\end{itemize}
	Вероятность мягкой посадки по модели YUCT составляет $<10^{-3}$.
	Критерий опровержения: если совокупное падение ВВП в 2026–2028 гг. составит $<5\%$, модель фальсифицирована.
	
	\section{Заключение и перспективы}
	
	Мы расширили YUCT на область $\Keff < 1$, ранее бывшую белым пятном теории. Ключевые достижения:
	
	\begin{itemize}
		\item \textbf{Отрицательная координация} определена как тахионная нестабильность с экспоненциальным
		ростом $\ln \Keff$.
		\item \textbf{Конструктивная нестабильность} различает абсолютный коллапс и возникновение нового порядка
		в зависимости от выживания фундаментальных словарей.
		\item Та же математика управляет фазовыми переходами (C2), биологической координацией (BCLM),
		эпигенетическим старением (BCLM‑EC), климатической волатильностью (Climat01) и экономическими циклами (F).
		\item \textbf{Хаос – это не противоположность координации; это режим низкого $\Keff$}.
		Следовательно, \textbf{управление $\Keff$ есть управление хаосом и энтропией} –
		принцип, имеющий непосредственное применение в биологии, материаловедении, климатической
		инженерии и экономике.
	\end{itemize}
	
	Приложение предоставляет всеобъемлющий набор фальсифицируемых предсказаний в разных дисциплинах,
	а также конкретные экспериментальные и наблюдательные проверки. Успешное подтверждение
	валидирует координационную парадигму и открывает дверь для \textbf{координированной инженерии хаоса}:
	намеренного проведения системы через конструктивную нестабильность для достижения более
	высокоупорядоченного состояния – универсальной стратегии для исцеления, омоложения,
	стабилизации климата и экономической устойчивости.
	
	\section*{Доступность данных и кода}
	
	Все таблицы данных, скрипты калибровки и прогнозные модели доступны по адресу
	\url{https://github.com/Alexey-Yakushev-YUCT/YPSDC}
	
	\section*{Благодарности}
	
	Автор благодарит участников семинара YUCT, разработчиков BCLM, BCLM‑EC, Climat01 и Приложения F,
	а также сообщества открытой науки за публикацию своих данных.
	
	\begin{thebibliography}{99}
		\bibitem{AppendixL} A. V. Yakushev, ``Appendix L: Fractal Coordination Error Scaling,'' YUCT, Zenodo (2026).
		\bibitem{AppendixY} A. V. Yakushev, ``Appendix Y: Mathematical Foundations of YUCT,'' YUCT, Zenodo (2026).
		\bibitem{AppendixG} A. V. Yakushev, ``Appendix G: Quantum Gravity Resolution,'' YUCT, Zenodo (2026).
		\bibitem{AppendixX} A. V. Yakushev, ``Appendix X: Generalisation of Shannon Theory,'' YUCT, Zenodo (2026).
		\bibitem{AppendixEhrenfest} A. V. Yakushev, ``Appendix Ehrenfest: Rotating Disk,'' YUCT, Zenodo (2026).
		\bibitem{AppendixJ} A. V. Yakushev, ``Appendix J: Half‑Integer Fermion Masses,'' YUCT, Zenodo (2026).
		\bibitem{AppendixT} A. V. Yakushev, ``Appendix T: Law of Evolution,'' YUCT, Zenodo (2026).
		\bibitem{AppendixC2} A. V. Yakushev, ``Appendix C2: Universal Scaling of Phase Transitions,'' YUCT, Zenodo (2026).
		\bibitem{AppendixBCLM} A. V. Yakushev, ``Appendix BCLM: Biological Coordination Lagrangian,'' YUCT, Zenodo (2026).
		\bibitem{AppendixBCLMEC} A. V. Yakushev, ``Appendix BCLM‑EC: Epigenetic Clocks,'' YUCT, Zenodo (2026).
		\bibitem{AppendixClimat01} A. V. Yakushev, ``Appendix Climat01: Coordination Climatology,'' YUCT, Zenodo (2026).
		\bibitem{AppendixF} A. V. Yakushev, ``Appendix F: Coordination Economics,'' YUCT, Zenodo (2026).
		\bibitem{Feigenbaum1978} M. J. Feigenbaum, \emph{J. Stat. Phys.} \textbf{19}, 25 (1978).
		\bibitem{Prigogine1977} I. Prigogine, G. Nicolis, \emph{Self‑Organization in Nonequilibrium Systems}. Wiley (1977).
		\bibitem{GISTEMP} GISTEMP Team, NASA GISS, \url{https://data.giss.nasa.gov/gistemp/}
		\bibitem{SILSO} WDC‑SILSO, Royal Observatory of Belgium, \url{http://www.sidc.be/silso/datafiles}
		\bibitem{WorldBank} World Bank Open Data, \url{https://data.worldbank.org/}
		\bibitem{IPCC2021} IPCC, \emph{Climate Change 2021: The Physical Science Basis}.
	\end{thebibliography}
	
\end{document}